\subsection{Matrix Algebra and Linear Systems}
Throughout this subsection, all matrices are over the same fixed field $\mathbb{F}$, and every size parameter is a positive integer.

\begin{definition} \label{def:matrices-and-matrix-operations} \hlt{Matrices and Matrix Operations}
\begin{enumerate}[label=(\roman*)]
\setlength{\itemsep}{0pt}
\item \hlt{Matrix space and equality}: $\operatorname{Mat}(m,n,\mathbb{F})\defeq\{[A_{j,k}]_{1\leq j\leq m,\ 1\leq k\leq n}:A_{j,k}\in\mathbb{F}\}$.\\
For $A,C$ in this space, $(A=C)\iff(\forall\,j\in\{1,\ldots,m\})(\forall\,k\in\{1,\ldots,n\})\,A_{j,k}=C_{j,k}$.
\item \hlt{Entries, rows, and columns}: if $A=[A_{j,k}]\in\operatorname{Mat}(m,n,\mathbb{F})$, then $A_{j,k}$ is its $(j,k)$-entry, $A_{j,\cdot}$ is row $j$, and $A_{\cdot,k}$ is column $k$.
\item \hlt{Transpose}: for $A\in\operatorname{Mat}(m,n,\mathbb{F})$, the transpose $A^t\in\operatorname{Mat}(n,m,\mathbb{F})$ satisfies $(\forall\,j\in\{1,\ldots,m\})(\forall\,k\in\{1,\ldots,n\})\,(A^t)_{k,j}\defeq A_{j,k}$.
\item \hlt{Main diagonal and diagonal matrix}: if $A=[A_{j,k}]\in\operatorname{Mat}(m,n,\mathbb{F})$ and $r=\min(m,n)$, its main diagonal is $(A_{1,1},\ldots,A_{r,r})$.\\
For $\lambda_1,\ldots,\lambda_n\in\mathbb{F}$, $\operatorname{diag}(\lambda_1,\ldots,\lambda_n)\in\operatorname{Mat}(n,n,\mathbb{F})$ has $(j,j)$-entry $\lambda_j$ and all off-diagonal entries zero.
\item \hlt{Symmetry}: a matrix $A\in\operatorname{Mat}(n,n,\mathbb{F})$ is symmetric iff $A^t=A$, and skew-symmetric iff $A^t=-A$.
\item \hlt{Triangularity}: a matrix $A\in\operatorname{Mat}(n,n,\mathbb{F})$ is upper triangular iff $(\forall\,j,k\in\{1,\ldots,n\})(j>k\Rightarrow A_{j,k}=0)$.\\
Lower triangular iff $(\forall\,j,k\in\{1,\ldots,n\})(j<k\Rightarrow A_{j,k}=0)$.\\
Diagonal iff $(\forall\,j,k\in\{1,\ldots,n\})(j\neq k\Rightarrow A_{j,k}=0)$.
\item \hlt{Linear operations}: for $A,C\in\operatorname{Mat}(m,n,\mathbb{F})$ and $a\in\mathbb{F}$, $(A+C)_{j,k}\defeq A_{j,k}+C_{j,k}$ and $(aA)_{j,k}\defeq aA_{j,k}$ for admissible $j,k$.
\item \hlt{Row-column pairing}: if $r_1,\ldots,r_n,c_1,\ldots,c_n\in\mathbb{F}$, $\mathbf{r}=(r_1,\ldots,r_n)$, and $\mathbf{c}=(c_1,\ldots,c_n)^t$, then $\mathbf{r}\mathbf{c}\defeq\sum\limits_{s=1}^{n}r_sc_s$.
\item \hlt{Matrix product}: for $A\in\operatorname{Mat}(m,n,\mathbb{F})$ and $B\in\operatorname{Mat}(n,p,\mathbb{F})$, the product $AB\in\operatorname{Mat}(m,p,\mathbb{F})$ satisfies\\
$(\forall\,j\in\{1,\ldots,m\})(\forall\,k\in\{1,\ldots,p\})\,(AB)_{j,k}\defeq\sum\limits_{s=1}^{n}A_{j,s}B_{s,k}$.
\item \hlt{Matrix-vector product}: identify $\mathbf{x}=(x_1,\ldots,x_n)^t\in\mathbb{F}^n$ with its column matrix.\\
For $A\in\operatorname{Mat}(m,n,\mathbb{F})$, define $A\mathbf{x}\in\mathbb{F}^m$ by $(\forall\,j\in\{1,\ldots,m\})\,(A\mathbf{x})_j\defeq\sum\limits_{k=1}^{n}A_{j,k}x_k$.
\item \hlt{Matrix units}: for $j\in\{1,\ldots,m\}$ and $k\in\{1,\ldots,n\}$, $E_{j,k}\in\operatorname{Mat}(m,n,\mathbb{F})$ has $(j,k)$-entry $1$ and all other entries zero.
\item \hlt{Identity and inverse}: $I_n\defeq\operatorname{diag}(1,\ldots,1)=[\delta_{j,k}]\in\operatorname{Mat}(n,n,\mathbb{F})$, where $\delta_{j,k}=1$ iff $j=k$. A matrix $B\in\operatorname{Mat}(n,n,\mathbb{F})$ is a two-sided inverse of $A\in\operatorname{Mat}(n,n,\mathbb{F})$ iff $AB=I_n=BA$; the matrix $A$ is invertible iff it has a two-sided inverse.
\end{enumerate}
\end{definition}

\begin{proposition} \label{prop:matrix-product-associativity-and-identities} \hlt{Basic Matrix Product Laws}\\
Let $A,A'\in\operatorname{Mat}(m,n,\mathbb{F})$, $B,B'\in\operatorname{Mat}(n,p,\mathbb{F})$, $C\in\operatorname{Mat}(p,q,\mathbb{F})$, and $a\in\mathbb{F}$. Then:
\begin{enumerate}[label=(\roman*)]
\setlength{\itemsep}{0pt}
\item $(AB)C=A(BC)$;
\item $I_mA=A=AI_n$;
\item $(A+A')B=AB+A'B$ and $A(B+B')=AB+AB'$;
\item $(aA)B=a(AB)=A(aB)$.
\end{enumerate}
\begin{proof}
\begin{enumerate}[label=(\roman*)]
\setlength{\itemsep}{0pt}
\item For admissible $j,k$, $((AB)C)_{j,k}=\sum\limits_{s=1}^{p}\sum\limits_{r=1}^{n}A_{j,r}B_{r,s}C_{s,k}=\sum\limits_{r=1}^{n}\sum\limits_{s=1}^{p}A_{j,r}B_{r,s}C_{s,k}=(A(BC))_{j,k}$ by finite distributivity.\\
Thus $(AB)C=A(BC)$ by entrywise equality.
\item For admissible $j,k$, $(I_mA)_{j,k}=\sum\limits_{r=1}^{m}\delta_{j,r}A_{r,k}=A_{j,k}=\sum\limits_{r=1}^{n}A_{j,r}\delta_{r,k}=(AI_n)_{j,k}$. Hence $I_mA=A=AI_n$.
\item For admissible $j,k$, finite distributivity implies $((A+A')B)_{j,k}=\sum\limits_{r=1}^{n}(A_{j,r}+A'_{j,r})B_{r,k}=(AB+A'B)_{j,k}$.\\
Also implies $(A(B+B'))_{j,k}=\sum\limits_{r=1}^{n}A_{j,r}(B_{r,k}+B'_{r,k})=(AB+AB')_{j,k}$. Entrywise equality proves both laws.
\item For admissible $j,k$, $((aA)B)_{j,k}=a\sum\limits_{r=1}^{n}A_{j,r}B_{r,k}=(a(AB))_{j,k}=(A(aB))_{j,k}$; hence the matrices are equal.
\end{enumerate}
\hlt{Structural Commentary}: Each law reduces to corresponding identities among finite scalar sums.\\
\hlt{Key Insight}: Entrywise equality transfers field arithmetic to compatible matrices.
\end{proof}
\end{proposition}

\begin{lemma} \label{lem:two-sided-inverse-uniqueness} \hlt{Uniqueness of Two-Sided Inverses}\\
Let $A,B,C\in\operatorname{Mat}(n,n,\mathbb{F})$. If $AB=I_n$ and $BC=I_n$, then $A=C$.
\begin{proof}
By Proposition \ref{prop:matrix-product-associativity-and-identities} and the hypotheses, $A=AI_n=A(BC)=(AB)C=I_nC=C$.\\
\hlt{Structural Commentary}: A common middle inverse forces the left and right inverses to coincide.\\
\hlt{Key Insight}: Associativity turns $AB=I_n=BC$ into $A=C$.
\end{proof}
\end{lemma}

\begin{definition} \hlt{Inverse Matrix Notation}\\
If $A\in\operatorname{Mat}(n,n,\mathbb{F})$ is invertible, its unique two-sided inverse by Lemma \ref{lem:two-sided-inverse-uniqueness} is denoted $A^{-1}$.\\
The set of invertible matrices in $\operatorname{Mat}(n,n,\mathbb{F})$ is denoted $\operatorname{Mat}(n,n,\mathbb{F})^{\times}$.
\end{definition}

\begin{proposition} \label{prop:matrix-spaces-are-vector-spaces} \hlt{Matrix Spaces}\\
For $m,n\in\mathbb{Z}_{>0}$, the following hold.
\begin{enumerate}[label=(\roman*)]
\setlength{\itemsep}{0pt}
\item The entrywise operations make $\operatorname{Mat}(m,n,\mathbb{F})$ a vector space over $\mathbb{F}$.
\item The family $(E_{j,k})_{1\leq j\leq m,\ 1\leq k\leq n}$ is a basis; hence $\dim(\operatorname{Mat}(m,n,\mathbb{F}))=mn$.
\end{enumerate}
\begin{proof}
\begin{enumerate}[label=(\roman*)]
\setlength{\itemsep}{0pt}
\item The zero matrix and entrywise negatives supply the additive identity and inverses.\\
Each remaining vector-space axiom holds in every entry by the corresponding field law, so it holds by matrix equality.
\item Every $A=[A_{j,k}]$ has the expansion $A=\sum\limits_{j=1}^{m}\sum\limits_{k=1}^{n}A_{j,k}E_{j,k}$. If $\sum\limits_{j=1}^{m}\sum\limits_{k=1}^{n}c_{j,k}E_{j,k}=0$, then its $(j,k)$-entry is $c_{j,k}=0$.\\
Hence the family spans and is linearly independent. Its $mn$ members form a basis, so the stated dimension is $mn$.
\end{enumerate}
\hlt{Structural Commentary}: Entrywise field laws establish the space, and matrix units establish its basis.\\
\hlt{Key Insight}: Each matrix unit isolates one independent scalar coordinate.
\end{proof}
\end{proposition}

\begin{proposition} \label{prop:matrix-product-entry-row-column} \hlt{Matrix Product Entries Are Row Times Column}\\
If $A\in\operatorname{Mat}(m,n,\mathbb{F})$ and $B\in\operatorname{Mat}(n,p,\mathbb{F})$, then $(\forall\,j\in\{1,\ldots,m\})(\forall\,k\in\{1,\ldots,p\})\,(AB)_{j,k}=A_{j,\cdot}B_{\cdot,k}$.
\begin{proof}
Both sides equal $\sum\limits_{r=1}^{n}A_{j,r}B_{r,k}$ by the definitions of matrix multiplication and row-column pairing.\\
\hlt{Structural Commentary}: Matrix multiplication is assembled entry-by-entry from row-column pairings.\\
\hlt{Key Insight}: Each entry of $AB$ sees one row of $A$ and one column of $B$.
\end{proof}
\end{proposition}

\begin{definition} \label{def:linear-map-induced-by-matrix} \hlt{Map Induced by a Matrix}\\
For $A\in\operatorname{Mat}(m,n,\mathbb{F})$, the \hlt{map induced by $A$} is $T_A:\mathbb{F}^n\to\mathbb{F}^m$ defined by $T_A(\mathbf{x})=A\mathbf{x}$.
\end{definition}

\begin{proposition} \label{prop:matrix-induced-map-linear} \hlt{Matrix-Induced Maps Are Linear}\\
If $A\in\operatorname{Mat}(m,n,\mathbb{F})$, then $T_A:\mathbb{F}^n\to\mathbb{F}^m$ is linear.
\begin{proof}
For $\mathbf{x},\mathbf{y}\in\mathbb{F}^n$, $a,b\in\mathbb{F}$, and $j\in\{1,\ldots,m\}$, $(T_A(a\mathbf{x}+b\mathbf{y}))_j=\sum\limits_{k=1}^{n}A_{j,k}(ax_k+by_k)=a(T_A\mathbf{x})_j+b(T_A\mathbf{y})_j$.\\
Coordinatewise equality proves $T_A(a\mathbf{x}+b\mathbf{y})=aT_A\mathbf{x}+bT_A\mathbf{y}$.\\
\hlt{Structural Commentary}: Matrix-vector multiplication preserves linear combinations.\\
\hlt{Key Insight}: Every matrix defines a canonical linear map between coordinate spaces.
\end{proof}
\end{proposition}

\begin{proposition} \label{prop:matrix-vector-product-column-combination} \hlt{Matrix-Vector Products as Column Combinations}\\
If $A=[\mathbf{a}_1\ \cdots\ \mathbf{a}_n]\in\operatorname{Mat}(m,n,\mathbb{F})$ and $\mathbf{x}=(x_1,\ldots,x_n)^t\in\mathbb{F}^n$, then $A\mathbf{x}=\sum\limits_{k=1}^{n}x_k\mathbf{a}_k$.
\begin{proof}
For $j\in\{1,\ldots,m\}$, both vectors have $j$th entry $\sum\limits_{k=1}^{n}A_{j,k}x_k$.\\
\hlt{Structural Commentary}: Multiplying by a column selects a linear combination of the columns of $A$.\\
\hlt{Key Insight}: The coordinates of $\mathbf{x}$ are precisely the column coefficients.
\end{proof}
\end{proposition}

\begin{proposition} \label{prop:matrix-product-by-columns} \hlt{Matrix Products by Columns}\\
If $A\in\operatorname{Mat}(m,r,\mathbb{F})$ and $B=[\mathbf{b}_1\ \cdots\ \mathbf{b}_n]\in\operatorname{Mat}(r,n,\mathbb{F})$, then $AB=[A\mathbf{b}_1\ \cdots\ A\mathbf{b}_n]$.
\begin{proof}
For $j\in\{1,\ldots,m\}$ and $k\in\{1,\ldots,n\}$, the $j$th entries of the $k$th column of $AB$ and of $A\mathbf{b}_k$ both equal $\sum\limits_{\ell=1}^{r}A_{j,\ell}B_{\ell,k}$.\\
\hlt{Structural Commentary}: Matrix products apply $A$ to columns of $B$.\\
\hlt{Key Insight}: The columns of $AB$ are obtained by applying $A$ to the columns of $B$.
\end{proof}
\end{proposition}

\begin{proposition} \label{prop:matrix-product-outer-product-decomposition} \hlt{Matrix Products as Sums of Outer Products}\\
Let $A=[\mathbf{a}_1\ \cdots\ \mathbf{a}_r]\in\operatorname{Mat}(m,r,\mathbb{F})$ and $B\in\operatorname{Mat}(r,n,\mathbb{F})$. If the $\ell$th row of $B$ is $\mathbf{b}_{\ell}^t$ with $\mathbf{b}_{\ell}\in\mathbb{F}^n$, then $AB=\sum\limits_{\ell=1}^{r}\mathbf{a}_{\ell}\mathbf{b}_{\ell}^t$.
\begin{proof}
For admissible $j,k$, the $(j,k)$-entry of $\mathbf{a}_{\ell}\mathbf{b}_{\ell}^t$ is $A_{j,\ell}B_{\ell,k}$, so the corresponding entry of the stated sum is $\sum\limits_{\ell=1}^{r}A_{j,\ell}B_{\ell,k}=(AB)_{j,k}$.\\
\hlt{Structural Commentary}: Column-row products decompose a matrix product into single column-row pieces.\\
\hlt{Key Insight}: Each shared index $\ell$ contributes one outer product.
\end{proof}
\end{proposition}

\begin{definition} \label{def:column-space-matrix-rank-nullity} \hlt{Column Space, Matrix Rank, and Nullity}\\
Let $A=[\mathbf{a}_1\ \cdots\ \mathbf{a}_n]\in\operatorname{Mat}(m,n,\mathbb{F})$.
\begin{enumerate}[label=(\roman*)]
\setlength{\itemsep}{0pt}
\item $\operatorname{col}(A)\defeq\operatorname{span}(\mathbf{a}_1,\ldots,\mathbf{a}_n)\subseteq\mathbb{F}^m$, and $\operatorname{row}(A)\defeq\operatorname{col}(A^t)\subseteq\mathbb{F}^n$.
\item $\ker(A)\defeq\{\mathbf{x}\in\mathbb{F}^n:A\mathbf{x}=\mathbf{0}\}$ is the \hlt{nullspace}, and $\operatorname{nullity}(A)\defeq\dim(\ker(A))$.
\item $\operatorname{rank}(A)\defeq\dim(\operatorname{col}(A))$ is the column rank, while $\dim(\operatorname{row}(A))$ is the row rank.
\item The matrix $A$ has \hlt{full rank} iff $\operatorname{rank}(A)=\min(m,n)$, \hlt{full column rank} iff $\operatorname{rank}(A)=n$, and \hlt{full row rank} iff $\operatorname{rank}(A)=m$.
\end{enumerate}
$\operatorname{range}(T_A)=\operatorname{col}(A)$ and $\operatorname{range}(T_{A^t})=\operatorname{row}(A)$ by Proposition \ref{prop:matrix-vector-product-column-combination}, while $\ker(T_A)=\ker(A)$ by definition.\\
These three sets are subspaces by Proposition \ref{prop:matrix-induced-map-linear} and Proposition \ref{prop:kernel-and-range-are-subspaces}, so the stated dimensions are defined.
\end{definition}

\begin{proposition} \label{prop:linear-system-solvability-and-affine-solution-set} \hlt{Solvability and Complete Solution Sets}\\
Let $A=[\mathbf{a}_1\ \cdots\ \mathbf{a}_n]\in\operatorname{Mat}(m,n,\mathbb{F})$ and $\mathbf{b}\in\mathbb{F}^m$. Then:
\begin{enumerate}[label=(\roman*)]
\setlength{\itemsep}{0pt}
\item The system $A\mathbf{x}=\mathbf{b}$ with $\mathbf{x}\in\mathbb{F}^n$ is solvable iff $\mathbf{b}\in\operatorname{col}(A)$.
\item If $A\mathbf{x}_p=\mathbf{b}$, its solution set is $\mathbf{x}_p+\ker(A)=\{\mathbf{x}_p+\mathbf{x}_0:\mathbf{x}_0\in\ker(A)\}$.
\end{enumerate}
\begin{proof}
\begin{enumerate}[label=(\roman*)]
\setlength{\itemsep}{0pt}
\item[$(i)$] ($\Rightarrow$) If $A\mathbf{x}=\mathbf{b}$ for $\mathbf{x}=(x_1,\ldots,x_n)^t$, then $\mathbf{b}=A\mathbf{x}=\sum\limits_{j=1}^{n}x_j\mathbf{a}_j\in\operatorname{col}(A)$ by Proposition \ref{prop:matrix-vector-product-column-combination}.
\item[$(i)$] ($\Leftarrow$) If $\mathbf{b}\in\operatorname{col}(A)$, then $(\exists\,x_1,\ldots,x_n\in\mathbb{F})\,\mathbf{b}=\sum\limits_{j=1}^{n}x_j\mathbf{a}_j$. For $\mathbf{x}=(x_1,\ldots,x_n)^t$, $A\mathbf{x}=\mathbf{b}$ by Proposition \ref{prop:matrix-vector-product-column-combination}.
\item[$(ii)$] If $A\mathbf{x}_p=\mathbf{b}$, then $(\forall\,\mathbf{x}\in\mathbb{F}^n)\,A\mathbf{x}=\mathbf{b}\iff A(\mathbf{x}-\mathbf{x}_p)=\mathbf{0}\iff\mathbf{x}-\mathbf{x}_p\in\ker(A)\iff\mathbf{x}\in\mathbf{x}_p+\ker(A)$ by Proposition \ref{prop:matrix-induced-map-linear}.
\end{enumerate}
\hlt{Structural Commentary}: Column expansion proves existence, while subtraction identifies the complete affine fiber.\\
\hlt{Key Insight}: Column space decides existence; kernel decides nonuniqueness.
\end{proof}
\end{proposition}

\begin{proposition} \label{prop:elementary-transpose-rules} \hlt{Elementary Transpose Rules}\\
Let $A,C\in\operatorname{Mat}(m,n,\mathbb{F})$ and $a\in\mathbb{F}$. Then:
\begin{enumerate}[label=(\roman*)]
\setlength{\itemsep}{0pt}
\item $(A^t)^t=A$;
\item $(A+C)^t=A^t+C^t$ and $(aA)^t=aA^t$;
\item $I_n^t=I_n$.
\end{enumerate}
\begin{proof}
\begin{enumerate}[label=(\roman*)]
\setlength{\itemsep}{0pt}
\item For admissible $j,k$, $((A^t)^t)_{j,k}=(A^t)_{k,j}=A_{j,k}$.
\item For admissible $j,k$, $((A+C)^t)_{j,k}=A_{k,j}+C_{k,j}=(A^t+C^t)_{j,k}$ and $((aA)^t)_{j,k}=aA_{k,j}=(aA^t)_{j,k}$.
\item For $j,k\in\{1,\ldots,n\}$, $(I_n^t)_{j,k}=\delta_{k,j}=\delta_{j,k}=(I_n)_{j,k}$.
\end{enumerate}
\hlt{Structural Commentary}: Each identity follows by swapping the two entry indices.\\
\hlt{Key Insight}: Transposition preserves entrywise operations and reverses index order.
\end{proof}
\end{proposition}

\begin{example} \label{ex:matrix-space-symmetric-skew-direct-sum} \hlt{Symmetric and Skew-Symmetric Decomposition}\\
Assume $2\neq0$ in $\mathbb{F}$ and set $M\defeq\operatorname{Mat}(n,n,\mathbb{F})$. Define $\operatorname{Sym}_n(\mathbb{F})\defeq\{A\in M:A^t=A\}$ and $\operatorname{Skew}_n(\mathbb{F})\defeq\{A\in M:A^t=-A\}$.\\
The map $\tau:M\to M$, $\tau(A)\defeq A^t$, is linear by Proposition \ref{prop:elementary-transpose-rules}(ii). Thus $\tau-I_M,\tau+I_M\in\LL(M)$ by Proposition \ref{prop:linear-maps-form-vector-space}.\\
Since $\operatorname{Sym}_n(\mathbb{F})=\ker(\tau-I_M)$ and $\operatorname{Skew}_n(\mathbb{F})=\ker(\tau+I_M)$, both are subspaces of $M$ by Proposition \ref{prop:kernel-and-range-are-subspaces}(i).\\
For $A\in M$, set $S_A\defeq\frac12(A+A^t)$ and $K_A\defeq\frac12(A-A^t)$. Then $S_A^t=S_A$, $K_A^t=-K_A$, and $A=S_A+K_A$ by Proposition \ref{prop:elementary-transpose-rules}.\\
Hence $M=\operatorname{Sym}_n(\mathbb{F})+\operatorname{Skew}_n(\mathbb{F})$. If $A\in\operatorname{Sym}_n(\mathbb{F})\cap\operatorname{Skew}_n(\mathbb{F})$, then $A=A^t=-A$ and $2A=0$, whence $A=0$ because $2\neq0$ is invertible in $\mathbb{F}$. Therefore $M=\operatorname{Sym}_n(\mathbb{F})\oplus\operatorname{Skew}_n(\mathbb{F})$ by Proposition \ref{prop:uniqueness-direct-sum-two-subspaces}.\\
\hlt{Structural Commentary}: Transpose splits a matrix into its even and odd parts.\\
\hlt{Key Insight}: The directness of the split uses the ability to divide by $2$.
\end{example}

\begin{proposition} \label{prop:transpose-reverses-products} \hlt{Transpose Reverses Products}\\
If $A\in\operatorname{Mat}(m,n,\mathbb{F})$ and $B\in\operatorname{Mat}(n,p,\mathbb{F})$, then $(AB)^t=B^tA^t$.
\begin{proof}
For $j\in\{1,\ldots,p\}$ and $k\in\{1,\ldots,m\}$, $((AB)^t)_{j,k}=(AB)_{k,j}=\sum\limits_{\ell=1}^{n}A_{k,\ell}B_{\ell,j}$.\\
The same sum equals $\sum\limits_{\ell=1}^{n}(B^t)_{j,\ell}(A^t)_{\ell,k}=(B^tA^t)_{j,k}$, so the matrices are equal.\\
\hlt{Structural Commentary}: Transposition swaps the order of composition encoded by matrix multiplication.\\
\hlt{Key Insight}: Rows become columns, so the product order reverses.
\end{proof}
\end{proposition}

\begin{corollary} \label{cor:matrix-inverse-product-and-transpose} \hlt{Matrix Inverses and Transposes Reverse Products}\\
Let $A,B\in\operatorname{Mat}(n,n,\mathbb{F})$ be invertible. Then:
\begin{enumerate}[label=(\roman*)]
\setlength{\itemsep}{0pt}
\item $AB$ is invertible and $(AB)^{-1}=B^{-1}A^{-1}$;
\item $A^t$ is invertible and $(A^t)^{-1}=(A^{-1})^t$.
\end{enumerate}
\begin{proof}
\begin{enumerate}[label=(\roman*)]
\setlength{\itemsep}{0pt}
\item For $D\defeq B^{-1}A^{-1}$, $D(AB)=I_n=(AB)D$ by Proposition \ref{prop:matrix-product-associativity-and-identities} and the inverse equations; hence $(AB)^{-1}=D$.
\item For $C\defeq(A^{-1})^t$, $CA^t=I_n=A^tC$ by Proposition \ref{prop:transpose-reverses-products} and Proposition \ref{prop:elementary-transpose-rules}(iii); hence $(A^t)^{-1}=C$.
\end{enumerate}
\hlt{Structural Commentary}: Each claim verifies the two inverse equations after reversing multiplication order.\\
\hlt{Key Insight}: Both inversion and transposition reverse products.
\end{proof}
\end{corollary}
